\documentclass{beamer}
\usetheme{metropolis}
% \setbeamertemplate{navigation symbols}{} 

\AtBeginEnvironment{alltt}{\setlength{\topsep}{0pt}}

\titlegraphic{\hfill\includegraphics[height=1cm]{figures/logo-dc.eps}}
\usepackage{textpos} 

\addtobeamertemplate{frametitle}{}{%
	\begin{textblock*}{0mm}(\textwidth-2cm,-1cm)
		\includegraphics[height=1cm]{figures/logofondonegro.eps}
    \end{textblock*}
}
\setbeamertemplate{section in toc}[ball unnumbered] % index with bullet points
\setbeamercolor{block body}{bg=structure!10} % visible blocks
\setbeamercolor{block title}{bg=structure!20}
\setbeamertemplate{blocks}[rounded][shadow] %round the edges of the blocks and cast a shadow
%-------------------------------------------------------
\usepackage[spanish]{babel} % Para separar correctamente las palabras
\usepackage[utf8]{inputenc} % Este paquete permite poner acentos y eñes usando codificación
\usepackage[most]{tcolorbox}
\usepackage{wrapfig}
\usepackage{multimedia}

\usepackage{multicol}
\usepackage{listings}
\lstset{
    language=C,
    backgroundcolor=\color{black!5}, % set backgroundcolor
    basicstyle=\footnotesize,% basic font setting
}

\usepackage{trimclip}
% Math

\newcommand{\card}[1]{\lvert#1\rvert}
\newcommand{\<}{\langle}
\renewcommand{\>}{\rangle}
\newcommand{\ldot}{\;.\;}
\newcommand{\lcom}{\,,\,}
\newcommand{\primeI}[1]{{#1}'}
\newcommand{\primeII}[1]{{#1}'\!\!\:'}
\newcommand{\into}{\mapsto}
\newcommand{\Nat}{\mathbb{N}}
\newcommand{\dotrightarrow}{\clipbox{0pt 0 {.675\width} 0}{\ensuremath{\rightarrow}} \makebox[1.1pt]{} {\cdot} {\cdot} {\cdot} \makebox[1.1pt]{} \clipbox{{.55\width} 0 0pt 0}{\ensuremath{\rightarrow}}}
\newcommand{\longdotrightarrow}{\clipbox{0pt 0 {.55\width} 0}{\ensuremath{\rightarrow}} \makebox[2pt]{} {\cdot} \makebox[1pt]{} {\cdot} \makebox[1pt]{} {\cdot} \makebox[2pt]{} \clipbox{{.45\width} 0 0pt 0}{\ensuremath{\rightarrow}}}

% Supervisory Control

\renewcommand{\l}{\ell}
\renewcommand{\ll}{{\bar{\l}}}
\renewcommand{\d}{\delta}
\newcommand{\init}[1]{\bar{#1}} % #1 state name
\newcommand{\state}[2]{{#1}^{#2}}
\newcommand{\s}[1]{\state{e}{#1}}

\newcommand{\D}{\rightarrow}
\newcommand{\B}{\rightsquigarrow}

\newcommand{\A}{\mathcal{A}}
\newcommand{\E}{\mathcal{E}}
\newcommand{\G}{\mathcal{G}}
\newcommand{\M}{\mathcal{M}}
\newcommand{\R}{\mathcal{R}}
\newcommand{\W}{\mathcal{W}}

\renewcommand{\L}{\mathcal{L}}
\renewcommand{\P}{\mathcal{P}}
\renewcommand{\S}{\mathcal{S}}

\newcommand{\w}{\ensuremath{\omega}}

\newcommand{\notstep}[2]{\overset{#1}{\not\D}_{#2}}
\newcommand{\step}[2]{\overset{#1}{\D}_{#2}}
\newcommand{\mstep}[2]{\overset{#1}{\twoheadrightarrow}_{#2}}
\newcommand{\lstep}[2]{\overset{#1}{\longrightarrow}_{#2}}
\newcommand{\walk}[2]{\overset{#1}{\Rightarrow}_{#2}}
\newcommand{\hop}[2]{\overset{#1}{\B}_{#2}}
\newcommand{\runw}[2]{\,\overset{#1}{\dotrightarrow}_{#2}\,}
\newcommand{\run}[3]{
\,\overset{#1\ldots#2}{\longdotrightarrow}_{#3}\,}
\newcommand{\runlst}[3]{
\overset{#1\quad#2}{\raisebox{.1pt}{-}{\,\cdot\cdot\cdot}\!\D}_{#3}}


\newcommand<>{\reveal}[1]{\mbox{}\visible#2{#1}} % ej: \reveal<5->{text} aparece de 5 en adelante

\title[DCS non-blocking]{Directed Controller Synthesis for Non-Maximal Blocking Requirements}
\author[Duran, Zanollo]{Matias Duran, Florencia Zanollo}
\date[12/5/21]{12 de Mayo de 2021}
%-------------------------------------------------------
% INTRO
% intro sobre qué queremos (programa que soluciona problemas por nosotros)
% qué toma de input (autómatas, modelo de la realidad, en general son varios)
% detalles de los autómatas (estado inicial, marcados, etc) con ejemplo de avión
% la composición explota, los algor clásicos rompen acá (i.e. decir que existe algor monolítico)
% idea general: componer mientras exploras y terminar al tener conclusión (con suerte antes de ver todo)
% [buscar] ejemplo simple para explicar composición (mostrar planta completa)

% DEFINICIONES (ANTECEDENTES)
% def pasos y corridas
% controlador
% vemos en particular problemas non-blocking (i.e. loop para ganadores, no-controlables no joden)
% estados ganadores y perdedores (estados desde donde hay una estrategia==controlador ganadora/perdedora)

% ON-THE-FLY
% explicar exploración parcial (lo visto hasta ahora)
% estados ganadores y perdedores ahora (antes era facil, ahora tenemos que ver a dónde van)
% el concepto de ``frontera'' (en el sentido de transiciones existentes sin explorar) top, bottom
% actualizamos estado solo cuando estamos seguros (ganador en bottom / perdedor en top)
% exploramos de a una transición [e -l-> e']
% esto actualiza antecesores de e
% idea central de que ganador necesita loop

% TESTS
% implementación en MTSA, explicación mini de la herramienta
% hicimos tests, medio TDD

% BENCHMARK
% usamos mismo benchmark que dany
% resultados vs DCS version 1.0 y la nuestra (DCS2)

% QUÉ VIMOS (CONCLUSIONES)
% entender el problema
% soluciones existentes, idea nueva
% nuevo algoritmo, ``demo'' de corrección y completitud
% implementación
% batería de tests
% benchmark confirma que mantiene eficiencia
% esta idea se puede aplicar a otro tipo de problemas (que no sean non-blocking), ej GR1 vuelvan prontos
%-------------------------------------------------------
\begin{document}
%-------------------------------------------------------
\begin{frame}[plain]

\titlepage

\end{frame}
%-------------------------------------------------------
\begin{frame}[plain]{Índice}
    \tableofcontents
\end{frame}
%-------------------------------------------------------
\section{Introducción}%-------------------------------------------------------
\begin{frame}{Titulo 1}
    contenido
\end{frame}
%-------------------------------------------------------

%-------------------------------------------------------
\section{Conocimiento previo}AI planning, control supervisado
nonblocking


\begin{definition}[Deterministic Automaton] \label{def:automaton}
	A \k{deterministic automaton} is a tuple $T = (S_T, A_T, \D_T, \init{t}, M_T)$, where:
	\begin{itemize*}[label=]
		
		\item $S_T$ is a \k{finite set of states};
		
		\item $A_T$ is the automaton \k{event set};
		
		\item $\D_T \C (S_T \x A_T \x S_T)$ is a \k{transition function};
		
		\item $\init{t} \in S_T$ is the \k{initial state}; and
		
		\item $M_T \C S_T$ is a set of \k{marked states}.
		
	\end{itemize*}
	
\end{definition}

\begin{notation}[Steps and Runs] \label{not:step}
	
	$\!\!$We denote $(t,\l,t') {\in}\!\! \D_T$ by $t \step{\l}{T} t'$ and call it a \k{step}.
	In turn, a \k{run} on a word $w = \l_0,\ldots,\l_k$ in $T$, is a sequence of steps such that $t_i \step{\l_i}{T} t_{i+1}$ for all $0 \leq i \leq k$, denoted by $t_0 \runw{w}{T} t_{k+1}$.
	
\end{notation}

\begin{definition} [Parallel Composition] \label{def:parcomp}
	The \k{parallel composition} $(\|)$ of two automata $T$ and $Q$ is an associative 
	symmetric operator that yields an automaton $T \| Q = (S_T {\times} S_Q, A_T {\cup} 
	A_Q, \D_{T\|Q}, \<\init{t},\init{q}\>, M_T {\times} M_Q)$, where $\D_{T\|Q}$ is the 
	smallest relation that satisfies the following rules (we omit the symmetric version of 
	the first rule):
	
	\begin{normalsize}
		\centering
		\vspace{-18pt}
		\hspace{-50pt}
		\begin{minipage}{0.30\linewidth}
			\[ 
			\frac{t \step{\l}{T} t'}{\<t,q\> {\step{\l}{T\|Q}} \<t'\!,\!q\> }{{\scriptstyle \l \in A_{T} {\setminus} A_{Q}}} 
			\]
		\end{minipage} 
		\hspace{40pt}
		%\begin{minipage}{0.30\linewidth}
		%\[ 
		%\frac{q \step{\l}{Q} q'}{\<t,q\> \step{\l}{T\|Q} \<t,q'\>}{{\scriptstyle \l \in A_{\!Q} 
		%{\setminus} A_{T}}} 
		%\]
		%\end{minipage} \\[-6pt]
		%\hspace{-20pt}
		\begin{minipage}{0.30\linewidth}
			\[ 
			\frac{t \step{\l}{T} t' \quad q \step{\l}{Q} q'}{\<t,q\> \step{\l}{T\|Q} \<t',q'\>}{{\scriptstyle \l \in A_{T} {\cap} A_{Q}}}
			\]
		\end{minipage} \\[15pt]
	\end{normalsize}
\end{definition}

\section{Controlador objetivo}

\begin{definition}[Supervisory Control Problem] \label{def:control-problem}
	A compositional \k{Supervisory Control Problem} is a tuple $\E = (E, A_E^C)$, where $E$ is a set of automata $\{E_0,\ldots,E_n\}$ (we may abuse notation and also use $E = (S_E,A_E,\D_E,\init{e},M_E)$ to refer to the composition $E_0\|\ldots\|E_n$), and $A_E^C \C A_E$ is the set of controllable events (i.e., $A_E^U = A_E \setminus A_E^C$ is the set of uncontrollable events).
	A solution for $\E$ is a supervisor $\sigma : A_E^* \into 2^{A_E}$, such that $\sigma$ is:
	\begin{itemize}[itemsep=4pt,topsep=-8pt]
		
		\item \k{Controllable}, namely $A_E^U \C \sigma(w)$ with $w \in A_E^*$; and
		
		\item \k{Nonblocking}, namely for every word $w \in \L^\sigma(E)$ there exists a non-empty word $w' \in A_E^*$ such that, the concatenation $ww' \in \L^\sigma(E)$ and $\init{e} \runw{\;ww'}{E} e_m$ with $e_m \in M_E$ (i.e., a marked state of $E$).
		
	\end{itemize}
	
\end{definition}


El problema a tratar consiste en encontrar para la planta de entrada, un controlador, es decir, un autómata con las siguientes características:

\begin{itemize}
	\item sub-autómata de la planta: todos los estados y transiciones del controlador existen en la planta compuesta.
	
	\item controlable: todas las transiciones no controlables de la planta se encuentran en el controlador.
	
	\item non-blocking: cada palabra válida para el controlador puede ser extendida por otra palabra no vacía para que su concatenación alcance un estado marcado.
\end{itemize}

Podemos pensar en un controlador non-blocking como un jugador optimista. Se encarga de no perder, y mientras tenga un futuro camino posible que lo lleva al destino buscado, considera que está ganando.

Es clave entender que en el problema a tratar, la posición de "tablas" del ajedrez, en la que ambos jugadores repiten sus jugadas 50 veces, se considera ganadora si todavía hay opción de dar un jaque mate. Si repetimos nuestras jugadas y todavía tengo dos torres considero que gané el partido, porque eventualmente mi oponente podría cansarse y dejarme ganar. Si repetimos nuestras jugadas pero solo tengo mi rey, no hay forma de dar mate, no puedo extender esta "palabra", esta partida, de forma de dar mate, y considero que perdí.

Es importante notar que como se busca que cualquier palabra sea extendible a otro estado marcado, lo que se busca es pasar por algún estado marcado infinitas veces. O sea, un estado 'e' marcado que tenga un camino para que el jugador pueda volver controlablemente al mismo estado 'e'.

Por esto, las estructuras claves que analizamos en nuestro algoritmo son los "loops", ya que los primeros estados ganadores son aquellos que están en un loop controlable con un estado marcado dentro. Luego señalizamos como ganadores también a cualquier estado que controlablemente alcanza un estado ganador.

También los "loops" son esenciales para encontrar los estados perdedores, ya que la única forma de que un estado sea perdedor es que no pueda alcanzar un estado ganador. En otras palabras, los estados perdedores son aquellos que forman parte de un loop que no tiene estados marcados ni transiciones salientes.

De forma más concreta, en nuestro algoritmo, un "loop" del cual el jugador no puede escapar, pero desde el cual existe un camino hacia un estado ganador, se considera ganador.


\subsection{Exploración on-the-fly}

El problema de síntesis de controlador ya tiene una solución clásica, por lo que la dificultad del trabajo no consistió en desarrollar un algoritmo que detectara estados ganadores y perdedores de un LTS totalmente explorado. 

\textcolor{red}{Explicar pto fijo clásico}

El conflicto reside en que al componer distintos DES, la cantidad de estados de la composición es exponencial respecto de los estados en los componentes. Esto es de suma relevancia ya que la solución clásica, que compone toda la planta para luego explorarla, tiene un límite de escalabilidad en el cual la composición de la planta llega al límite de tiempo o memoria, y nunca se llega a la exploración.

Para combatir esto, la exploración on-the-fly clasifica estados como ganadores o perdedores durante la composición. Se espera que con esto sea posible, en primer lugar, cortar la exploración de una rama de la planta que ya se sabe que es perdedora o ganadora, reduciendo así la memoria y tiempo necesarios. Pero más aún, si el estado inicial fuera marcado como ganador o perdedor antes de la composición completa de la planta, ni siquiera sería necesario completar el proceso de composición.

Para incrementar las ramas podadas se utiliza una heurística de exploración Best First Search \textcolor{red}{[REF Dani]} que busca ganar controlablemente o perder no controlablemente, para garantizar con la menor exploración posible que el estado actual es ganador o perdedor.

En el peor caso, se perdió tiempo en los puntos fijos, intentando clasificar estados, y se realiza una última vez el algoritmo clásico con la planta totalmente explorada. Esto garantiza la completitud del algoritmo, como se detalla en mayor profundidad en el capítulo~\ref{chpt:paper}.

















%-------------------------------------------------------
\section{On-the-fly}% ON-THE-FLY
% explicar exploración parcial (lo visto hasta ahora)
% estados ganadores y perdedores ahora (antes era fácil, ahora tenemos que ver a dónde van)
% el concepto de ``frontera'' (en el sentido de transiciones existentes sin explorar) top, bottom
% actualizamos estado solo cuando estamos seguros (ganador en bottom / perdedor en top)
% exploramos de a una transición [e -l-> e']
% esto actualiza antecesores de e
% idea central de que ganador necesita loop
%-------------------------------------------------------
\begin{frame}{Exploración parcial}
    i.e. es lo visto hasta el momento
\end{frame}
%-------------------------------------------------------
\begin{frame}{Estados ganadores y perdedores (bis)}
    estados ganadores y perdedores ahora (antes era fácil, ahora tenemos que ver a dónde van)
\end{frame}
%-------------------------------------------------------
\begin{frame}{Frontera optimista y pesimista}
    def el concepto de ``frontera'' (en el sentido de transiciones existentes sin explorar) top, bottom
\end{frame}
%-------------------------------------------------------
\begin{frame}{Conclusión sobre un estado}
    actualizamos estado solo cuando estamos seguros (ganador en bottom / perdedor en top)
\end{frame}
%-------------------------------------------------------
\begin{frame}{Detalles de la exploración}
    exploramos de a una transición [e -l-> e']. 
    ejemplo de la tesis como animación?. 
\end{frame}
%-------------------------------------------------------
\begin{frame}{Sólo actualizamos antecesores}
    explicar por qué o lema o ej de imágen, ver la mejor manera
\end{frame}
%-------------------------------------------------------
\begin{frame}{Ganador necesita tener un loop}
    idea central de que ganador necesita loop, esto no sería mejor decirlo antes? por ej cuando hablamos de ganadores?
\end{frame}
%-------------------------------------------------------

%-------------------------------------------------------
\section{Implementación y Tests}% TESTS
% implementación en MTSA, explicación mini de la herramienta
% hicimos tests, medio TDD
%-------------------------------------------------------
\begin{frame}{Herramienta MTSA}
    \begin{itemize}
    	\item El LaFHIS, donde hicimos esta tesis, trabaja hace años con MTSA, una herramienta propia para resolver problemas de control con LTS.
    	\item Implementamos el algoritmo en java, agregándolo a las capacidades de MTSA. Esto permitió ejecutar el benchmark de la próxima sección.
    \end{itemize}
\end{frame}
%-------------------------------------------------------
\begin{frame}{Test Driven Development}
	\begin{itemize}
		\item Para ganar seguridad en nuestro código, y encontrar errores, fuimos armando una batería de tests.
		\item Con cada error encontrado en la implementación o el algoritmo, armábamos un nuevo test que detectara ese error, y luego lo arreglábamos
		\item Quedaron 49 tests pequeños diseñados a mano para correr rápido y presentar las condiciones más problemáticas para nuestro algoritmo
		\item los tests quedan en la herramienta como tests de regresión para prevernir problemas por futuros cambios
	\end{itemize}
\end{frame}
%-------------------------------------------------------
\begin{frame}{Ejemplo de MTSA en uso}
	\begin{figure}
		\hspace{-2cm}
		\includegraphics[width=\textwidth]{figures/HPWindow-ej.png}
	\end{figure}
	\pause
	\begin{picture}(0,0)
		\put(70,-10){
			\includegraphics[width=\textwidth]{figures/controller-drawing-ej.png}
		}
	\end{picture}
\end{frame}
%-------------------------------------------------------
\begin{frame}{Animando controladores}
	Video de controlador animado?
\end{frame}
%-------------------------------------------------------

%-------------------------------------------------------
\section{Benchmark}\begin{frame}{Benchmark}
    \begin{itemize}
     \item Para realizar las pruebas de performance usamos el mismo conjunto de problemas utilizado para medir la versión anterior del algoritmo de exploración, recompilados por Daniel Ciolek en su tesis doctoral.
     
     \item Es un conjunto de seis tipos de problemas bastante clásicos, cada uno con dos partes parametrizables (n, k), en función de observar hasta qué tamaño de problema (n*k) soporta el algoritmo.
     
     \item DCS y DCS2 representan la versión anterior/nueva del algoritmo respectivamente.
     
     \item Utilizamos distintas estrategias a la hora de explorar, algunas más complejas (MA, RA) desarrolladas por Ciolek y una simple (BFS). La última fue agregada para obtener una idea de cuánto se puede mejorar cambiando la forma de explorar.
    \end{itemize}

\end{frame}
%-------------------------------------------------------
\begin{frame}{Performance (TL, DP, CM)}
    \begin{figure}
        \includegraphics[width=\textwidth]{figures/benchmark1.png}
    \end{figure}

    Transfer Line, Dinning Philosophers, Cat and Mouse
    
\end{frame}
%-------------------------------------------------------
\begin{frame}{Performance (AT, BW, TA)}
    \begin{figure}
        \includegraphics[width=\textwidth]{figures/benchmark2.png}
    \end{figure}

    Air-Traffic Management, Bidding Workflow, Travel Agency
\end{frame}

%-------------------------------------------------------
\section{Conclusión}
\begin{frame}[plain]{Qué vimos}
    \begin{itemize}[<+->]
    \item Definición del problema
    \item Soluciones existentes y la idea ``nueva'' (exploración on-the-fly)
    \item Nuevo algoritmo, parte de su demostración (corrección y completitud)
    \item Implementación en MTSA
    \item Batería de tests, TDD
    \item Benchmark y resultados versus la versión anterior. No se perdió eficiencia, teniendo en cuenta la confianza ganada en correctitud. En la tesis se encuentran además resultados de benchmark versus otras herramientas del estado del arte.
    \end{itemize}

    \begin{block}<7->{}
    La idea de exploración on-the-fly, y gran parte de la estructura del algoritmo, se puede aplicar a otro tipo de problemas, ej: GR1 (próximamente).
    \end{block}
\end{frame}
%-------------------------------------------------------
\begin{frame}[plain]{\null} % el null esta para forzar el header, alto trucazo
    \begin{columns}
    \begin{column}{0.5\textwidth}
        \Huge ¿Preguntas?\\
        
        \normalsize Gracias, vuelvan prontos\\ 
        \scriptsize (para ver GR1)
    \end{column}
    \begin{column}{0.5\textwidth}
        \includegraphics[scale=0.5]{figures/apu.png}
    \end{column}
    \end{columns}
\end{frame}
%-------------------------------------------------------
\begin{frame}{Bonus track: Detalles de la exploración}    
    \begin{columns}
    \begin{column}{0.8\textwidth}
        \begin{figure}
        \centering
            \begin{overprint}
                \onslide<1>\includegraphics[scale=0.55]{../figures/ejemplo_on-the-fly/5.pdf}
                \onslide<2>\includegraphics[scale=0.55]{../figures/ejemplo_on-the-fly/1.pdf}
                \onslide<3>\includegraphics[scale=0.55]{../figures/ejemplo_on-the-fly/2.pdf}
                \onslide<4>\includegraphics[scale=0.55]{../figures/ejemplo_on-the-fly/3.pdf}
                \onslide<5->\includegraphics[scale=0.55]{../figures/ejemplo_on-the-fly/4.pdf}
            \end{overprint}
        \end{figure}
    \end{column}
    \begin{column}{0.4\textwidth}
        \hspace{-50pt}
        \begin{overprint}
            \onslide<2>Empieza por el estado \textbf{0} y mira los pasos:
                \begin{itemize}
                 \item $0 \step{1}{} 1$
                 \item $1 \step{2}{} 2$
                 \item $0 \step{3(u)}{} 3$
                 \item $1 \step{4(c)}{} 3$
                \end{itemize}

            \onslide<3>Explora el paso $3 \step{5}{} 1$\\
            Cierra loop pero todavía no hay estado marcado (u objetivo).
                 
            \onslide<4>Mira los pasos:
                \begin{itemize}
                 \item $2 \step{6}{} 4$
                 \item $4 \step{7}{} 2$
                \end{itemize}
                Cierra loop que incluye el estado marcado \textbf{4}, encontramos una parte ganadora, propagamos esa información.
                
            \onslide<5->Explora $3 \step{8}{} 5$, encuentra que \textbf{5} es perdedor y propaga la información. Llega al inicial entonces concluye que no hay controlador.
        \end{overprint}
    \end{column}
    \end{columns}

\end{frame}
%-------------------------------------------------------
\end{document}
